% !TeX program = lualatex
% !TeX encoding = utf8
% !TeX spellcheck = uk_UA
% !BIB program = biber

\documentclass[a4paper,12pt]{article}

\usepackage{fontsetup}
\usepackage[english, russian, ukrainian]{babel}


\usepackage{amsmath}
\usepackage{graphicx}
\usepackage{geometry}
\geometry{top=2cm, bottom=2cm, left=2cm, right=2cm}
\usepackage{hyperref}

\title{Конспект: Нестійкості в плазмі}
\author{Фізика плазми}
\date{}

\begin{document}

\maketitle

\section{Гідродинамічне та квазігідродинамічне наближення}

\begin{itemize}
    \item \textbf{Гідродинамічне наближення}: опис рівноваги та руху газів за допомогою сили тиску. Застосовне до \textbf{густої плазми}, де часті зіткнення забезпечують передачу тиску.
    \item \textbf{Квазігідродинамічне наближення}: опис \textbf{розрідженої плазми} за допомогою сили тиску поперек магнітного поля. Фізичний механізм передачі тиску тут інший: дрейфовий струм та струм намагнічування.
    \item Умови застосування: ті самі, що для дрейфового або адіабатичного наближення.
\end{itemize}

\subsection{Кількісний механізм передачі тиску}

Якщо поперечний тиск плазми $p_{\perp}$ і напруженість магнітного поля $H$ змінюються в напрямку $x$ (перпендикулярному до поля), то перпендикулярно до $H$ та $x$ тече струм:
\[
|j| = |j_H + j_M| = c \left[ \frac{p_{\perp}}{H^2} \cdot \frac{dH}{dx} + \frac{d}{dx} \left( \frac{p_{\perp}}{H} \right) \right]
\]
Після спрощення:
\[
|j| = \frac{c}{H} \cdot \frac{dp_{\perp}}{dx}
\]
Порівнюючи з формулою дрейфового струму $|j_D| = cn \frac{F_{\perp}}{H}$, отримуємо ефективну силу:
\[
F_{\perp} = -\frac{1}{n} \cdot \frac{dp_{\perp}}{dx}
\]
Таким чином, дрейфовий струм і струм намагнічування еквівалентні дії сили тиску.

\section{Аналогія з гідродинамічною нестійкістю (важка рідина над легкою)}

\begin{figure}[h!]
    \centering
    % \includegraphics[width=0.4\textwidth]{two_layers.png}
    \caption{Двошарова система (а) та її нестійкість (б). Важка рідина (сірчана кислота) зверху, легка (вода) знизу.}
    \label{fig:two_layers}
\end{figure}

\begin{itemize}
    \item Система перебуває в рівновазі (тиск нижнього шару врівноважує вагу верхнього), але це рівновага \textbf{нестійка}.
    \item Енергетично вигідніше, коли важка рідина знизу, а легка зверху.
    \item Випадкове мале збурення (крапля важкої рідини опускається) запускає \textbf{конвекцію} – перемішування шарів.
    \item Приклад: осел Буридана – вибір між двома однаковими оберемками сіна; мала флуктуація визначає вибір.
\end{itemize}

\section{Магнітогідродинамічна (МГД) нестійкість}

\subsection{Нестійкість на межі плазма–магнітне поле (жолобкова нестійкість)}

\begin{itemize}
    \item \textbf{Умова виникнення}: межа між провідною плазмою (всередині якої немає магнітного поля) та зовнішнім магнітним полем.
    \item \textbf{Скін-шар}: тонкий шар, де магнітне поле різко змінюється від нуля до значення у зовнішньому просторі.
    \item \textbf{Енергетична невигідність}: магнітні силові лінії, випуклі в бік плазми (центр кривизни всередині плазми), розтягнуті. Обмін місцями плазми та поля зменшує енергію магнітного поля.
    \item \textbf{Механізм}: плазма є діамагнітною (виштовхується в область слабшого поля). Якщо поле зменшується з відстанню від плазми – виникає нестійкість.
    \item \textbf{Характер}: конвективна нестійкість, аналогічна нестійкості важкої рідини над легкою, але замість сили тяжіння – магнітне натягнення силових ліній.
\end{itemize}

\subsection{Нестійкість пінча}

\begin{itemize}
    \item Лінійний пінч (струм тече вздовж плазмового циліндра, створюючи азимутальне магнітне поле) є нестійким.
    \item \textbf{Деформації}:
    \begin{itemize}
        \item $m = 0$: \textbf{«шийка»} (перетяжка) – симетричне звуження та розширення.
        \item $m = 1$: \textbf{«змійка»} – вигин циліндра.
        \item $m \ge 2$: вищі азимутальні моди.
    \end{itemize}
    \item Механізм: магнітний тиск зовні (азимутальне поле) та всередині (подовжнє поле, якщо є) створюють сили, що підсилюють деформацію.
\end{itemize}

\section{Стабілізація вмороженим магнітним полем}

\subsection{Умови стабілізації}

\begin{itemize}
    \item \textbf{Вморожене поле}: поле, «приклеєне» до плазми через високу провідність.
    \item \textbf{Механізм стабілізації}: будь-який рух плазми поперек силових ліній згинає їх, що вимагає енергії (пружність силових ліній). Якщо кінці силових ліній закріплені на твердих провідниках (електродах), деформація утруднена.
    \item Аналогія: силові лінії – як арматура в залізобетоні.
\end{itemize}

\subsection{Стабілізація лінійного пінча}

\begin{itemize}
    \item Найкраща стабілізація: всередині плазми – \textbf{лише подовжнє} поле, струм тече по \textbf{поверхні} (просторово розділені поля).
    \item Для стабілізації мод $m \neq 1$ потрібне достатньо сильне подовжнє поле.
    \item Для моди $m = 0$ (шийка): $H_{\text{внутр}} \ge \frac{1}{\sqrt{2}} H_{\text{зовн, азим}} \approx 0.707 H_{\text{зовн}}$.
    \item Моди $m \ge 2$ стабілізуються легше.
\end{itemize}

\subsection{Обмеження стабілізації}

\begin{itemize}
    \item У слабоіонізованій плазмі можливе зворотне: магнітне поле може спричиняти нестійкість струму (пондеромоторна сила).
    \item Навіть в ідеально провідній плазмі вморожене поле не дає повної стабілізації для всіх збурень.
\end{itemize}

\section{Жолобкова нестійкість та способи боротьби з нею}

\subsection{Природа жолобкової нестійкості}

\begin{itemize}
    \item \textbf{Перестановочна нестійкість}: плазма та магнітні силові лінії обмінюються місцями. Силова лінія переміщується як ціле, зберігаючи форму.
    \item Збурення межі мають вигляд \textbf{жолобків}, витягнутих вздовж силових ліній.
    \item Механізм: плазма ніби «розсуває» силові лінії та «просочується» у проміжки між ними.
\end{itemize}

\subsection{Способи стабілізації}

\begin{enumerate}
    \item \textbf{Закріплення кінців силових ліній}: вморожування в тверді провідники (металеві стінки або електроди). Тоді перестановка неможлива.
    \item \textbf{Створення стабілізуючого поля всередині плазми}, спрямованого під кутом до зовнішнього (найкраще – \textbf{перпендикулярно}). На межі утворюється «плетена сітка» (корзинка), яка не дає плазмі просідати (рис. 27а).
\end{enumerate}

\subsection{Роль розмиття межі}

\begin{itemize}
    \item Для ідеально провідної плазми межа різка. Реальна плазма має скінченну провідність, тому межа розмита дифузією поля.
    \item \textbf{Парадоксальний ефект}:
    \begin{itemize}
        \item Якщо \textbf{немає} стабілізуючого поля всередині – розмиття межі \textbf{стабілізує} (згладжує градієнт).
        \item Якщо \textbf{є} перпендикулярне стабілізуюче поле – розмиття призводить до \textbf{перемішування полів} (винтове поле, рис. 27б), що відновлює жолобкову нестійкість.
    \end{itemize}
    \item \textbf{Висновок}: вплив факторів залежить від конкретних умов; критерії стійкості є складними та делікатними.
\end{itemize}

\section{Адiабатичні пастки (магнітні дзеркала) та жолобкова нестійкість}

\begin{itemize}
    \item У магнітній пастці з пробками (дзеркалами) в перехідній області силові лінії неминуче \textbf{випуклі назовні}.
    \item У цій області «язики» плазми просочуються між жолобками (витягнутими вздовж силових ліній), що веде до втрати плазми на стінку.
    \item Досліди показують: плазма в адіабатичній пастці розпадається досить швидко – це пояснюється жолобковою нестійкістю.
    \item \textbf{Методи боротьби}:
    \begin{itemize}
        \item «Приморожування» кінців силових ліній до металевих провідників.
        \item Заміна адіабатичної пастки на \textbf{пастку із зустрічними полями} (квадрупольна пастка). Там поле всюди ввігнуте, тому жолобкова нестійкість відсутня (але є «дірки»: дві на полюсах і кільцева вздовж екватора).
    \end{itemize}
\end{itemize}

\section{Дифузія протилежних полів}

\begin{itemize}
    \item При контакті областей з \textbf{протилежно направленими} магнітними полями (однакової величини) їх взаємна дифузія призводить до утворення \textbf{нейтрального шару}, де магнітне поле дорівнює нулю.
    \item Магнітний тиск у нейтральному шарі дорівнює нулю, тоді як з обох боків він додатний.
    \item Це явище виникає при дії змінних полів, у сонячній атмосфері, а також пропонувалося як метод стабілізації пінча (штучне створення зовнішнього поля, протилежного внутрішньому).
\end{itemize}

\section*{Висновки}

\begin{enumerate}
    \item У плазмі існує широкий клас \textbf{гідромагнітних нестійкостей}, аналогічних гідродинамічним (важка рідина над легкою).
    \item \textbf{Вморожене магнітне поле} може стабілізувати плазму, діючи як арматура, але не усуває всі типи збурень.
    \item \textbf{Жолобкова (перестановочна) нестійкість} є ключовою для адіабатичних пасток; боротьба з нею потребує спеціальної геометрії полів або закріплення силових ліній.
    \item Вплив \textbf{розмиття межі} (скінченна провідність) є неоднозначним і залежить від наявності стабілізуючого поля – це ілюструє складність критеріїв стійкості.
    \item \textbf{Дифузія протилежних полів} створює нейтральні шари, що важливо для фізики Сонця та стабілізації пінчів.
\end{enumerate}

\end{document}